\section{Desarrollo y preguntas abiertas}
\subsection{Cómo ha evolucionado el núcleo}
Comencé por fijar el nombre, el alcance exclusivamente escrito, dieciséis
signos, una gramática inicial y diez ejemplos. Al practicar, Colección y Fuente
me resultaron más difíciles de escribir, aunque recordaba y entendía los signos
por separado.

Después amplié el núcleo con seis funciones, reglas de composición, catálogos,
veinte enunciados de prueba y un cuaderno. Adopté las formas simplificadas de
Colección y Fuente porque me resultaron más cómodas y distinguibles. También
ajusté las uniones de Mientras y Espacio.

\subsection{Qué quiero mejorar con la práctica}
\begin{itemize}
\item Comprobar si reconozco las formas dentro de expresiones completas.
\item Identificar cualquier confusión entre los signos.
\item Simplificar las expresiones que me resulten demasiado largas al escribir.
\item Decidir qué composiciones frecuentes me conviene abreviar con una definición.
\item Ampliar el vocabulario a partir de ideas que realmente quiera expresar.
\end{itemize}

\subsection{Mi criterio para avanzar}
Busco que cada revisión resuelva una dificultad observada y conserve el
significado de las expresiones anteriores. Uso T01--T20 como controles para
comparar cambios. Las comprobaciones técnicas aportan resultados concretos;
mi comodidad y comprensión requieren práctica personal.

Mantengo abiertos el vocabulario general, otras clases de preguntas, una
expresión temporal más rica, argumentos extensos, métodos de entrada y tipografía.
Por ahora sigo concentrándome exclusivamente en la escritura.

\subsection{Después de consolidar el manual}
Ya tengo reunidos los signos, el trazado, las reglas y los ejemplos, con el
historial en anexos. Mi siguiente paso es practicar la lectura de expresiones
completas y retomar M01 otro día. Quiero usar esa experiencia para orientar
las próximas mejoras sin esperar a cubrir todas las disciplinas.
